\usepackage{luacode}

\begin{luacode*}
  results = {}

  -- Split CSV line while preserving empty fields.
  local function split_csv(line)
  local fields = {}
  for field in (line .. ","):gmatch("([^,]*),") do
  fields[#fields + 1] = field
  end
  return fields
  end

  -- Normalize a row key: integers stay as integer strings, everything else stays raw.
  local function normalize_key(s)
  local n = tonumber(s)
  if n and n == math.floor(n) then
  return string.format("%d", n)
    elseif n then
    return tostring(n)
    else
    return s
    end
    end

    -- Load a pivoted CSV where rows are timepoints (or pseudonyms) and columns are modality combinations.
    -- Expected structure:
    -- row 1: ehr,<flags...>
    -- row 2: waves,<flags...>
    -- row 3: thermal,<flags...>
    -- row 4: time_min,,,,
    -- row 5+: <time>,<value_for_combo_1>,...,<value_for_combo_n>
    function load_pivot_csv(dataset, condition, filepath)
    if not results[dataset] then results[dataset] = {} end
    if not results[dataset][condition] then results[dataset][condition] = {} end

    local ehr_flags = {}
    local waves_flags = {}
    local thermal_flags = {}
    local combo_keys = {}

    local row_idx = 0
    for line in io.lines(filepath) do
    row_idx = row_idx + 1
    local fields = split_csv(line)

    if row_idx == 1 then
    for col = 2, #fields do
    ehr_flags[col] = fields[col]
    end
    elseif row_idx == 2 then
    for col = 2, #fields do
    waves_flags[col] = fields[col]
    end
    elseif row_idx == 3 then
    for col = 2, #fields do
    thermal_flags[col] = fields[col]
    combo_keys[col] = thermal_flags[col] .. "-" .. ehr_flags[col] .. "-" .. waves_flags[col]
    if not results[dataset][condition][combo_keys[col]] then
    results[dataset][condition][combo_keys[col]] = {}
    end
    end
    elseif row_idx >= 5 then
    local row_key = normalize_key(fields[1])
    if row_key and row_key ~= "" then
    for col = 2, #fields do
    local key = combo_keys[col]
    local val = tonumber(fields[col])
    if key then
    if val then
    results[dataset][condition][key][row_key] = string.format("%.2f", val)
      else
      results[dataset][condition][key][row_key] = "??"
      end
      end
      end
      end
      end
      end
      end

      -- Load the new by-time AUROC results.
      load_pivot_csv("Test_mean", "auroc_by_time", "new_results/auroc_epoch2_by_time_mean.csv")
      load_pivot_csv("Test_std", "auroc_by_time", "new_results/auroc_epoch2_by_time_std.csv")
      load_pivot_csv("Validation_mean", "auroc_by_time", "new_results/auroc_epoch2_by_time_mean.csv")
      load_pivot_csv("Validation_std", "auroc_by_time", "new_results/auroc_epoch2_by_time_std.csv")

      -- Load the ranks snapshot at 90min.
      load_pivot_csv("Test_mean", "ranks_90min", "new_results/ranks_epoch2_snapshot_90min_mean.csv")
      load_pivot_csv("Test_std", "ranks_90min", "new_results/ranks_epoch2_snapshot_90min_std.csv")
      load_pivot_csv("Validation_mean", "ranks_90min", "new_results/ranks_epoch2_snapshot_90min_mean.csv")
      load_pivot_csv("Validation_std", "ranks_90min", "new_results/ranks_epoch2_snapshot_90min_std.csv")

      -- Load a two-method comparison CSV where:
      -- row 1: <row_label>,Our,Vanilla VAE
      -- row 2+: <key>,<our_val>,<vanilla_val>
      -- Stored as results[dataset][condition][method][row_key].
      function load_comparison_csv(dataset, condition, filepath)
      if not results[dataset] then results[dataset] = {} end
      if not results[dataset][condition] then results[dataset][condition] = {} end

      local method_keys = {}
      local row_idx = 0
      for line in io.lines(filepath) do
      row_idx = row_idx + 1
      local fields = split_csv(line)
      if row_idx == 1 then
      for col = 2, #fields do
      method_keys[col] = fields[col]
      if not results[dataset][condition][method_keys[col]] then
      results[dataset][condition][method_keys[col]] = {}
      end
      end
      else
      local row_key = normalize_key(fields[1])
      if row_key and row_key ~= "" then
      for col = 2, #fields do
      local key = method_keys[col]
      local val = tonumber(fields[col])
      if key then
      if val then
      results[dataset][condition][key][row_key] = string.format("%.2f", val)
        else
        results[dataset][condition][key][row_key] = "??"
        end
        end
        end
        end
        end
        end
        end

        -- Our vs. Vanilla VAE comparison (all modalities active).
        load_comparison_csv("Test_mean", "auroc_comparison", "new_results/auroc_comparison_mean.csv")
        load_comparison_csv("Test_std", "auroc_comparison", "new_results/auroc_comparison_std.csv")
        load_comparison_csv("Test_mean", "ranks_comparison_90min", "new_results/ranks_comparison_90min_mean.csv")
        load_comparison_csv("Test_std", "ranks_comparison_90min", "new_results/ranks_comparison_90min_std.csv")

        -- Generate comparison ranks table rows from patient_list.
        -- Each row: label & \gradient{our_mean}{our_mean±std} & \gradient{vanilla_mean}{vanilla_mean±std} \\
        function generate_comparison_ranks_rows()
        local methods = { "Our", "Vanilla VAE" }
        for _, p in ipairs(patient_list) do
        local pseudo = p.pseudo
        local label = p.label
        local cells = {}
        for _, method in ipairs(methods) do
        local m = results["Test_mean"] and results["Test_mean"]["ranks_comparison_90min"]
        and results["Test_mean"]["ranks_comparison_90min"][method]
        and results["Test_mean"]["ranks_comparison_90min"][method][pseudo] or "??"
        local s = results["Test_std"] and results["Test_std"]["ranks_comparison_90min"]
        and results["Test_std"]["ranks_comparison_90min"][method]
        and results["Test_std"]["ranks_comparison_90min"][method][pseudo]
        local ms = m
        if m ~= "??" and s then ms = m .. " \\pm " .. s end
        cells[#cells + 1] = "\\gradient{" .. m .. "}{" .. ms .. "}"
        end
        tex.print("        " .. label .. " &")
        tex.print("        " .. cells[1] .. " &")
        tex.print("        " .. cells[2] .. " \\\\")
        end
        end

        -- Load patient descriptions: pseudonym -> label, description
        patient_labels = {}
        patient_descriptions = {}
        patient_list = {}
        do
        local first = true
        for line in io.lines("new_results/patient_descriptions.csv") do
        if first then
        first = false
        else
        local label = line:match("^([^,]*)")
        local rest = line:sub(#label + 2)
        local pseudo = rest:match("^([^,]*)")
        local desc = rest:sub(#pseudo + 2)
        patient_labels[pseudo] = label
        patient_descriptions[pseudo] = desc
        patient_list[#patient_list + 1] = { pseudo = pseudo, label = label, desc = desc }
        end
        end
        end

        -- Generate ranks table rows from patient_list.
        -- Each row: label & \gradient{mean}{mean±std} & ... \\
        function generate_ranks_rows()
        local combos = {
          {0,1,0}, {0,1,1}, {0,0,1}, {1,0,0}, {1,1,0}, {1,0,1}, {1,1,1}
        }
        for i, p in ipairs(patient_list) do
        local pseudo = p.pseudo
        local label = p.label
        local cells = {}
        for _, c in ipairs(combos) do
        local key = c[1] .. "-" .. c[2] .. "-" .. c[3]
        local m = results["Test_mean"] and results["Test_mean"]["ranks_90min"]
        and results["Test_mean"]["ranks_90min"][key]
        and results["Test_mean"]["ranks_90min"][key][pseudo] or "??"
        local s = results["Test_std"] and results["Test_std"]["ranks_90min"]
        and results["Test_std"]["ranks_90min"][key]
        and results["Test_std"]["ranks_90min"][key][pseudo]
        local ms = m
        if m ~= "??" and s then ms = m .. " \\pm " .. s end
        cells[#cells + 1] = "\\gradient{" .. m .. "}{" .. ms .. "}"
        end
        tex.print("        " .. label .. " &")
        for j = 1, #cells - 1 do
        tex.print("        " .. cells[j] .. " &")
        end
        tex.print("        " .. cells[#cells] .. " \\\\")
        end
        end

        -- Generate patient description table rows from patient_list.
        -- Each row: label & description \\
        function generate_patient_rows()
        for _, p in ipairs(patient_list) do
        tex.print("        " .. p.label .. " & " .. p.desc .. " \\\\")
        end
        end

      \end{luacode*}

      \newcommand{\patientlabel}[1]{%
        \directlua{
          local pseudo = "#1"
          local label = patient_labels[pseudo]
          if label then
          tex.print(label)
          else
          tex.print("??")
          end
        }%
      }

      \newcommand{\patientdesc}[1]{%
        \directlua{
          local pseudo = "#1"
          local desc = patient_descriptions[pseudo]
          if desc then
          tex.print(desc)
          else
          tex.print("??")
          end
        }%
      }

      \newcommand{\rankstablerows}{\directlua{generate_ranks_rows()}}
      \newcommand{\patienttablerows}{\directlua{generate_patient_rows()}}
      \newcommand{\generatecompranksrows}{\directlua{generate_comparison_ranks_rows()}}

      \newcommand{\getcm}[4]{%
        \directlua{
          local base = "#1"
          local condition = "#2"
          local method = "#3"

          local function normalize_time(t)
          local n = tonumber(t)
          if n and n == math.floor(n) then return tostring(math.floor(n))
          elseif n then return tostring(n)
          else return t end
          end

          local row_key = normalize_time("#4")
          local mean_dataset = base .. "_mean"

          local mean_val = results[mean_dataset]
          and results[mean_dataset][condition]
          and results[mean_dataset][condition][method]
          and results[mean_dataset][condition][method][row_key]

          if mean_val then
          tex.print(mean_val)
          else
          tex.print("??")
          end
        }%
      }

      \newcommand{\getcms}[4]{%
        \directlua{
          local base = "#1"
          local condition = "#2"
          local method = "#3"

          local function normalize_time(t)
          local n = tonumber(t)
          if n and n == math.floor(n) then return tostring(math.floor(n))
          elseif n then return tostring(n)
          else return t end
          end

          local row_key = normalize_time("#4")

          local mean_dataset = base .. "_mean"
          local std_dataset = base .. "_std"

          local mean_val = results[mean_dataset]
          and results[mean_dataset][condition]
          and results[mean_dataset][condition][method]
          and results[mean_dataset][condition][method][row_key]

          local std_val = results[std_dataset]
          and results[std_dataset][condition]
          and results[std_dataset][condition][method]
          and results[std_dataset][condition][method][row_key]

          if mean_val and std_val then
          tex.print(mean_val .. " ± " .. std_val)
          else
          tex.print("??")
          end
        }%
      }

      \newcommand{\getms}[6]{%
        \directlua{
          local base = "#1"
          local condition = "#2"
          local key = "#3-#4-#5"

          local function normalize_time(t)
          local n = tonumber(t)
          if n and n == math.floor(n) then return tostring(math.floor(n))
          elseif n then return tostring(n)
          else return t end
          end

          local time = normalize_time("#6")

          local mean_dataset = base .. "_mean"
          local std_dataset = base .. "_std"

          local mean_val = results[mean_dataset]
          and results[mean_dataset][condition]
          and results[mean_dataset][condition][key]
          and results[mean_dataset][condition][key][time]

          local std_val = results[std_dataset]
          and results[std_dataset][condition]
          and results[std_dataset][condition][key]
          and results[std_dataset][condition][key][time]

          if mean_val and std_val then
          tex.print(mean_val .. " ± " .. std_val)
          else
          tex.print("??")
          end
        }%
      }

      \newcommand{\getm}[6]{%
        \directlua{
          local base = "#1"
          local condition = "#2"
          local key = "#3-#4-#5"

          local function normalize_time(t)
          local n = tonumber(t)
          if n and n == math.floor(n) then return tostring(math.floor(n))
          elseif n then return tostring(n)
          else return t end
          end

          local time = normalize_time("#6")
          local mean_dataset = base .. "_mean"

          local mean_val = results[mean_dataset]
          and results[mean_dataset][condition]
          and results[mean_dataset][condition][key]
          and results[mean_dataset][condition][key][time]

          if mean_val then
          tex.print(mean_val)
          else
          tex.print("??")
          end
        }%
      }

      \newcommand{\getstd}[6]{%
        \directlua{
          local base = "#1"
          local condition = "#2"
          local key = "#3-#4-#5"

          local function normalize_time(t)
          local n = tonumber(t)
          if n and n == math.floor(n) then return tostring(math.floor(n))
          elseif n then return tostring(n)
          else return t end
          end

          local time = normalize_time("#6")
          local std_dataset = base .. "_std"

          local std_val = results[std_dataset]
          and results[std_dataset][condition]
          and results[std_dataset][condition][key]
          and results[std_dataset][condition][key][time]

          if std_val then
          tex.print(std_val)
          else
          tex.print("??")
          end
        }%
      }